% संपूर्ण मराठी ग्रंथातील प्रकरण 43: विन्यासमीमांसा आणि चिन्हार्थमीमांसा
% हा संपादनयोग्य स्रोतखंड आहे. संकलनासाठी संपूर्ण स्रोतसंग्रहातील
% अक्षररूपे, पूर्वभाग, चिन्हव्याख्या आणि आकृती-स्रोत आवश्यक आहेत.
% मूळ: Open Logic Project; मजकूर आणि रूपांतर CC BY 4.0.
% यंत्रानुवाद, दुरुस्ती आणि तपासणी: OpenAI Codex —
% GPT-5.6 Sol आणि GPT-6 Sol, दोन्ही Ultra effort.
% दुरुस्ती-पुनरावलोकन: GPT-6.1 Sol, Ultra effort.
\chapter{विन्यासमीमांसा आणि चिन्हार्थमीमांसा}\label{nml:syn:chap}
\section{परिचय}\label{nml:syn:int:sec}
मोडल तर्कशास्त्र \emph{मोडल विधाने} आणि त्यांच्यातील तार्किक
निष्पन्नता-संबंध यांचा अभ्यास करते. मोडल विधानांची पुढील
उदाहरणे पाहा:
\begin{enumerate}
\item $2+2=4$ हे अनिवार्य आहे.
\item उद्या पाऊस पडणे अनिवार्यपणे संभाव्य आहे.
\item जर~$\varphi$ अनिवार्यपणे संभाव्य असेल, तर~$\varphi$ संभाव्य आहे.
\end{enumerate}
संभाव्यता आणि अनिवार्यता या एकमेव मोडलता नाहीत: इतर एकस्थानी
संयोजकही मोडलता मानले जातात. उदाहरणार्थ, ``$\varphi$ असे असायला हवे,''
``पुढे~$\varphi$ घडेल,'' ``Dana ला~$\varphi$ माहीत आहे,'' किंवा
``Dana चा~$\varphi$ वर विश्वास आहे.''
मोडल तर्कशास्त्र प्रथम अरिस्टॉटल यांच्या \emph{De
  Interpretatione} मध्ये दिसते. अनिवार्य असल्यास संभाव्य असते,
पण उलट नाही; संभाव्यता आणि अनिवार्यता एकमेकींच्या साहाय्याने
परिभाषित करता येतात; $\varphi \land \psi$ संभाव्यतः सत्य असेल तर
$\varphi$ संभाव्यतः सत्य असते आणि $\psi$ संभाव्यतः सत्य असते,
पण उलट नाही; तसेच $\varphi \to \psi$ अनिवार्य असेल, तर $\varphi$
अनिवार्य असताना $\psi$ ही अनिवार्य असते, या गोष्टी प्रथम
त्यांच्या लक्षात आल्या.
मोडल तर्कशास्त्राकडे पाहण्याचा पहिला आधुनिक दृष्टिकोन
C.~I. Lewis यांच्या कार्यातून पुढे आला. Lewis आणि Langford
यांच्या \emph{Symbolic Logic} (1932) या ग्रंथात या कार्याची
परिणती झाली. वास्तविक अभिव्यंजन वापरून अभिव्यंजनाचे प्रतिनिधित्व
करण्याबाबत Lewis आणि Langford समाधानी नव्हते: $\varphi \lif \psi$
हे ``$\varphi$ पासून $\psi$ निष्पन्न होते'' याचे अपुरे प्रतिनिधित्व आहे.
त्याऐवजी ``अनिवार्यतः, जर $\varphi$ तर $\psi$'' असे अभिव्यंजनाचे
लक्षण मांडून त्यांनी ते $\varphi \strictif \psi$ असे दर्शवले.
या विविध गुणधर्मांची मांडणी करताना Lewis यांनी पाच भिन्न
मोडल प्रणाली ओळखल्या: \Log{S1}, \ldots, \Log{S4}, \Log{S5};
त्यांतील शेवटच्या दोन अजूनही वापरात आहेत.
Lewis आणि Langford यांचा दृष्टिकोन पूर्णतः \emph{विन्यासमीमांसक}
होता: त्यांनी योग्य वाटणारी स्वयंसिद्धके व नियम निवडले आणि
त्यांच्या साहाय्याने काय सिद्ध करता येते याचा अभ्यास केला.
चिन्हार्थमीमांसक दृष्टिकोन बराच काळ उपलब्ध नव्हता.
नंतर Rudolf Carnap यांनी \emph{Meaning and Necessity} (1947)
मध्ये \emph{अवस्था-वर्णन} ही संकल्पना वापरून पहिला प्रयत्न केला.
अवस्था-वर्णन म्हणजे त्या वर्णनात ``सत्य'' असलेल्या आण्विक
वाक्यांचा संग्रह. कोणत्याही वाक्य~$\varphi$ साठी सत्यतेची व्याख्या
विस्तारित केल्यानंतर, $\varphi$ सर्व अवस्था-वर्णनांत सत्य असेल
तर ते \emph{अनिवार्यतः सत्य} आहे, अशी Carnap यांची व्याख्या आहे.
मात्र त्यांचा दृष्टिकोन \emph{पुनरावृत्त} मोडलता हाताळू शकत
नव्हता: ``संभाव्यतः अनिवार्यतः \ldots संभाव्यतः $\varphi$''
या रूपाची वाक्ये नेहमी सर्वांत आतील मोडलतेपुरती सीमित होत.
मोडल चिन्हार्थमीमांसेतील मोठी प्रगती Saul Kripke यांच्या
``A Completeness Theorem in Modal Logic'' (JSL 1959) या
लेखामुळे झाली. एखादे विधान ``सर्व संभाव्य जगांत'' सत्य
असेल तर ते अनिवार्यतः सत्य असते, या लाइब्निझ यांच्या कल्पनेवर
Kripke यांनी आपले कार्य आधारले. परंतु त्या साध्या अर्थलक्षणातही
Carnap यांच्या दृष्टिकोनासारखीच अडचण आहे: जग~$w$ येथे एखादे
विधान अनिवार्यतः सत्य आहे की नाही, हे $w$ च्या निवडीवर
(किंवा अवस्था-वर्णन~$s$ वर) मुळीच अवलंबून नसते.
म्हणून Kripke यांनी जगे \emph{प्राप्यता संबंध}~$R$ ने
संबंधित आहेत असे मानले; आणि ``अनिवार्यतः $\varphi$'' या रूपाचे
विधान जग~$w$ येथे तेव्हा आणि केवळ तेव्हाच सत्य असते,
जेव्हा $w$ पासून \emph{प्राप्य असलेल्या} प्रत्येक जग~$w'$
येथे $\varphi$ सत्य असते. या दृष्टिकोनाचे कोणतेही रूप वापरणाऱ्या
चिन्हार्थमीमांसेला क्रिप्के चिन्हार्थमीमांसा म्हणतात;
तिच्यामुळे मोडल तर्कशास्त्रांचा, अनेकवचनात, झपाट्याने विकास
होऊ शकला.
क्रिप्के चिन्हार्थमीमांसेनुसार अर्थ लावल्यास, मोडल तर्कशास्त्र
\emph{संबंधी रचना} ``आतून'' कशा दिसतात हे दाखवते. संबंधी
रचना म्हणजे द्विस्थानी संबंध असलेला संच; उदाहरणार्थ,
वर्गातील विद्यार्थ्यांचा सामाजिक सुरक्षा क्रमांकानुसार
लावलेला संच ही संबंधी रचना आहे. पण संबंधी रचना अनेक
क्षेत्रांत आढळतात. जगाच्या अवस्थांमधील सापेक्ष संभाव्यतेखेरीज,
एखाद्या व्यक्तीच्या ज्ञानाच्या अवस्था ज्ञानविषयक संभाव्यतेने
संबंधित असू शकतात किंवा गतिशील प्रणालीच्या अवस्थांमध्ये
स्थित्यंतरे असू शकतात. या सर्वांचे प्रतिरूपण मोडल तर्कशास्त्राने
करता येते: पहिल्या प्रकारातून नेहमीचे सत्यविषयक मोडल तर्कशास्त्र,
तर इतर प्रकारांतून ज्ञानविषयक तर्कशास्त्र, गतिशील तर्कशास्त्र
इत्यादी मिळतात.
मोडल तर्कशास्त्रज्ञ ``अनुरूपता सिद्धांत'' म्हणतात त्या एका
विशिष्ट दृष्टीकोनावर आपण लक्ष केंद्रित करतो. Kripke यांच्या
आरंभीच्या महत्त्वाच्या निष्कर्षांपैकी एक असा की प्राप्यता
संबंध~$R$ चे अनेक गुणधर्म, उदाहरणार्थ संक्रमणशीलता किंवा
सममिती, योग्य ``मोडल रूपबंधां''च्या साहाय्याने
\emph{मोडल भाषेतच} व्यक्त करता येतात. उदाहरणार्थ,
$R$ ची परावर्तिता ``जर अनिवार्यतः~$\varphi$, तर~$\varphi$'' या
रूपबंधाशी ``अनुरूप'' आहे, असे मोडल तर्कशास्त्रज्ञ म्हणतात.
\Ax{D}, \Ax{T}, \Ax{B}, \Ax{4} आणि~\Ax{5} हे रूपबंध
एकत्र करून मिळणाऱ्या अभिजात मोडल तर्कशास्त्रांपैकी
काहींच्या, उदाहरणार्थ \Log{S4} आणि \Log{S5}, अनुरूपता
सिद्धांताचा आपण मुख्यतः अभ्यास करू.
\section{मूलभूत मोडल तर्कशास्त्राची भाषा}\label{nml:syn:lan:sec}
\begin{defn}
मोडल तर्कशास्त्राच्या मूलभूत भाषेत पुढील गोष्टी असतात:
\begin{enumerate}
  \item असत्यता साठीचा विधानीय स्थिरांक~$\lfalse$.
  \item सत्यता साठीचा विधानीय स्थिरांक~$\ltrue$.
  \item विधानीय चलांचा गणनीय अनंत संच: $\Obj
    p_0$, $\Obj p_1$, $\Obj p_2$, \dots
  \item विधानीय संयोजक: \startycommalist
  \ycomma $\lnot$ (निषेधन)
  \ycomma $\land$ (संधी)
  \ycomma $\lor$ (विकल्प)
  \ycomma $\lif$ (शर्त)
  \ycomma $\liff$ (द्विशर्त).
  \item मोडल परिकर्मी $\Box$.
  \item मोडल परिकर्मी $\Diamond$.
\end{enumerate}
\end{defn}
\begin{defn}
मूलभूत मोडल भाषेतील \emph{सूत्रे} पुढीलप्रमाणे
  विगमनाने परिभाषित होतात:
\begin{enumerate}
\item $\lfalse$ हे आण्विक सूत्र आहे.
\item $\ltrue$ हे आण्विक सूत्र आहे.
\item प्रत्येक विधानीय चल $\Obj p_i$ हे (आण्विक) सूत्र आहे.
\item $\varphi$ हे सूत्र असेल, तर $\lnot \varphi$
  हे सूत्र आहे.
\item $\varphi$ आणि $\psi$ ही सूत्रे असतील, तर $(\varphi \land
  \psi)$ हे सूत्र आहे.
\item $\varphi$ आणि $\psi$ ही सूत्रे असतील, तर $(\varphi \lor \psi)$
  हे सूत्र आहे.
\item $\varphi$ आणि $\psi$ ही सूत्रे असतील, तर $(\varphi \lif \psi)$
  हे सूत्र आहे.
\item $\varphi$ आणि $\psi$ ही सूत्रे असतील, तर $(\varphi \liff \psi)$
  हे सूत्र आहे.
\item $\varphi$ हे सूत्र असेल, तर $\Box \varphi$
  हे सूत्र आहे.
\item $\varphi$ हे सूत्र असेल, तर $\Diamond \varphi$
  हे सूत्र आहे.
\item यांशिवाय दुसरे काहीही सूत्र नाही.
\end{enumerate}
\end{defn}
सूत्र~$\varphi$ मध्ये $\Box$ किंवा $\Diamond$ नसेल,
तर त्याला \emph{मोडलतारहित} म्हणतो.
\section{एकाच वेळी केलेले आदेशन}\label{nml:syn:sub:sec}
सूत्र~$\varphi$ मधील विधानीय चल च्या सर्व
आस्थितींच्या जागी दुसरे सूत्र घालून मिळणारे सूत्र हा
$\varphi$ चा \emph{दृष्टांत} असतो. वैधता आणि निष्पन्नता
यांवर चर्चा करताना सूत्रांच्या दृष्टांतांचा वारंवार
उल्लेख करावा लागेल. म्हणून ही संकल्पना नेमकी परिभाषित
करणे उपयुक्त आहे.
\begin{defn}\label{nml:syn:sub:def:subst-inst}
  $\varphi$ हे असे मोडल सूत्र असू द्या की त्यातील सर्व
  विधानीय चले ही $p_1$, \dots, $p_n$ यांपैकीच
  आहेत; तसेच $\theta_1$, \dots, $\theta_n$ हीदेखील मोडल सूत्रे
  असू द्या. मग
  $\SSubst{\varphi}{\subst{\theta_1}{p_1}, \dots, \subst{\theta_n}{p_n}}$
  म्हणजे $\varphi$ मध्ये प्रत्येक $p_i$ च्या जागी $\theta_i$ चे
  एकाच वेळी आदेशन करून मिळणारे सूत्र, अशी व्याख्या करतो.
  औपचारिकरीत्या ही $\varphi$ वरील विगमनाने केलेली व्याख्या आहे:
  \begin{enumerate}
    \item \indcase{\varphi}{\lfalse}{
        $\SSubst{\indfrm}{\subst{\theta_1}{p_1}, \dots,
          \subst{\theta_n}{p_n}}$ हे $\lfalse$ असते.}
    \item \indcase{\varphi}{\ltrue}{
        $\SSubst{\indfrm}{\subst{\theta_1}{p_1}, \dots,
          \subst{\theta_n}{p_n}}$ हे $\ltrue$ असते.}
    \item \indcase{\varphi}{q}{$\SSubst{\indfrm}{\subst{\theta_1}{p_1}, \dots,
        \subst{\theta_n}{p_n}}$ हे $q$ असते, जर $i
      = 1$, \dots,~$n$ साठी $q \not\ident p_i$ असेल.}
    \item \indcase{\varphi}{p_i}{$\SSubst{\indfrm}{\subst{\theta_1}{p_1},
        \dots, \subst{\theta_n}{p_n}}$ हे $\theta_i$ असते.}
    \item \indcase{\varphi}{\lnot \psi}{
        $\SSubst{\indfrm}{\subst{\theta_1}{p_1}, \dots, \subst{\theta_n}{p_n}}$ हे
        $\lnot \SSubst{\psi}{\subst{\theta_1}{p_1}, \dots, \subst{\theta_n}{p_n}}$ असते.}
    \item \indcase{\varphi}{(\psi \land
        \chi)}{$\SSubst{\indfrm}{\subst{\theta_1}{p_1}, \dots,
          \subst{\theta_n}{p_n}}$ हे पुढील सूत्र असते:
        \[(\SSubst{\psi}{\subst{\theta_1}{p_1}, \dots, \subst{\theta_n}{p_n}} \land
          \SSubst{\chi}{\subst{\theta_1}{p_1}, \dots, \subst{\theta_n}{p_n}}).\]}
    \item \indcase{\varphi}{(\psi \lor
          \chi)}{$\SSubst{\indfrm}{\subst{\theta_1}{p_1}, \dots,
          \subst{\theta_n}{p_n}}$ हे पुढील सूत्र असते:
        \[(\SSubst{\psi}{\subst{\theta_1}{p_1}, \dots, \subst{\theta_n}{p_n}} \lor
          \SSubst{\chi}{\subst{\theta_1}{p_1}, \dots, \subst{\theta_n}{p_n}}).\]}
    \item \indcase{\varphi}{(\psi \lif
          \chi)}{$\SSubst{\indfrm}{\subst{\theta_1}{p_1}, \dots,
          \subst{\theta_n}{p_n}}$ हे पुढील सूत्र असते:
        \[(\SSubst{\psi}{\subst{\theta_1}{p_1}, \dots, \subst{\theta_n}{p_n}} \lif
          \SSubst{\chi}{\subst{\theta_1}{p_1}, \dots, \subst{\theta_n}{p_n}}).\]}
    \item \indcase{\varphi}{(\psi \liff
          \chi)}{$\SSubst{\indfrm}{\subst{\theta_1}{p_1}, \dots,
          \subst{\theta_n}{p_n}}$ हे पुढील सूत्र असते:
        \[(\SSubst{\psi}{\subst{\theta_1}{p_1}, \dots, \subst{\theta_n}{p_n}} \liff
          \SSubst{\chi}{\subst{\theta_1}{p_1}, \dots, \subst{\theta_n}{p_n}}).\]}
    \item \indcase{\varphi}{\Box \psi}{
        $\SSubst{\indfrm}{\subst{\theta_1}{p_1}, \dots, \subst{\theta_n}{p_n}}$ म्हणजे
        $\Box
        \SSubst{\psi}{\subst{\theta_1}{p_1}, \dots, \subst{\theta_n}{p_n}}.$}
    \item \indcase{\varphi}{\Diamond \psi}{
        $\SSubst{\indfrm}{\subst{\theta_1}{p_1}, \dots, \subst{\theta_n}{p_n}}$ म्हणजे
        $\Diamond
        \SSubst{\psi}{\subst{\theta_1}{p_1}, \dots, \subst{\theta_n}{p_n}}.$}
  \end{enumerate}
  सूत्र $\SSubst{\varphi}{\subst{\theta_1}{p_1}, \dots,
    \subst{\theta_n}{p_n}}$ याला $\varphi$ चा \emph{आदेशन-दृष्टांत}
  म्हणतात.
\end{defn}

\begin{ex}
  समजा $\varphi$ हे $p_1 \lif \Box(p_1 \land p_2)$,
  $\theta_1$ हे $\Diamond(p_2 \lif p_3)$ आणि $\theta_2$ हे
  $\lnot\Box p_1$ आहे. मग
  $\SSubst{\varphi}{\subst{\theta_1}{p_1}, \subst{\theta_2}{p_2}}$ म्हणजे
  \begin{align*}
    \Diamond(p_2 \lif p_3) &
    \lif \Box(\Diamond(p_2 \lif p_3) \land \lnot\Box p_1)
    \intertext{तर $\SSubst{\varphi}{\subst{\theta_2}{p_1}, \subst{\theta_1}{p_2}}$ म्हणजे}
    \lnot\Box p_1 &
    \lif \Box(\lnot\Box p_1 \land \Diamond(p_2 \lif p_3))
    \intertext{लक्षात घ्या: एकाच वेळी केलेले आदेशन
      सर्वसाधारणपणे क्रमाने केलेल्या आदेशनासारखे नसते.
      उदाहरणार्थ, वरील
      $\SSubst{\varphi}{\subst{\theta_1}{p_1}, \subst{\theta_2}{p_2}}$ ची
      $\Subst{(\Subst{\varphi}{\theta_1}{p_1})}{\theta_2}{p_2}$ शी तुलना करा; नंतरचे असे आहे:}
    & \Subst{(\Diamond(p_2 \lif p_3) \lif
      \Box(\Diamond(p_2 \lif p_3) \land p_2))}{\lnot\Box p_1}{p_2},
      \text{ म्हणजे}\\
    \Diamond(\lnot\Box p_1 \lif p_3) &
    \lif \Box(\Diamond(\lnot\Box p_1
    \lif p_3) \land \lnot\Box p_1)\\
    \intertext{आणि $\Subst{(\Subst{\varphi}{\theta_2}{p_2})}{\theta_1}{p_1}$ शीही तुलना करा:}
    & \Subst{(p_1 \lif \Box(p_1 \land \lnot\Box p_1))}
      {\Diamond(p_2 \lif p_3)}{p_1}, \text{ म्हणजे}\\
    \Diamond(p_2 \lif p_3) &
    \lif \Box(\Diamond(p_2
    \lif p_3) \land \lnot\Box\Diamond(p_2 \lif p_3)).
  \end{align*}
\end{ex}












\section{संबंधी प्रतिमाने}\label{nml:syn:rel:sec}

सामान्य मोडल तर्कशास्त्रांच्या चिन्हार्थमीमांसेची मूलभूत संकल्पना
म्हणजे \emph{संबंधी प्रतिमान}. त्यात द्विस्थानी ``प्राप्यता संबंधाने''
संबंधित जगांचा संच आणि कोणते विधानीय चले कोणत्या
जगांत ``सत्य'' मानायची हे ठरवणारे मूल्यांकन असते.

\begin{defn}
  मूलभूत मोडल भाषेचे \emph{प्रतिमान} म्हणजे त्रिकूट $\mModel{M}
  = \tuple{W, R, V}$; येथे
  \begin{enumerate}
  \item $W$ हा ``जगांचा'' रिकामा नसलेला संच आहे,
  \item $R$ हा~$W$ वरील द्विस्थानी प्राप्यता संबंध आहे, आणि
  \item $V$ हे प्रत्येक विधानीय चल~$p$ ला संभाव्य जगांचा संच~$V(p)$ नेमून देणारे फलन आहे.
  \end{enumerate}
  $Rww'$ असेल तेव्हा $w'$ हे~$w$ पासून \emph{प्राप्य}
  आहे असे म्हणतो. $w \in V(p)$ असेल तेव्हा $p$ हे~$w$ येथे
  \emph{सत्य} आहे असे म्हणतो.
\end{defn}

संबंधी चिन्हार्थमीमांसेचा मोठा फायदा असा की प्रतिमाने
\ref{nml:syn:rel:fig:simple} मधील साध्या आकृतीसारख्या आकृत्यांनी दाखवता
येतात. जगांसाठी गाठी वापरतात; $w$ पासून~$w'$ कडे बाण
असेल तेव्हा आणि केवळ तेव्हाच जग~$w'$ हे~$w$ पासून प्राप्य
असते. शिवाय $w \in V(p)$ असेल तर गाठीला (जगाला)~$\mTrue{p}$,
अन्यथा~\mFalse{p} अशी खूण करतो. \ref{nml:syn:rel:fig:simple} मध्ये
$W = \{w_1, w_2, w_3\}$, $R = \{\tuple{w_1, w_2},\tuple{w_1,
  w_3}\}$, $V(p) = \{w_1, w_2\}$ आणि $V(q) = \{w_2\}$
असलेले प्रतिमान दाखवले आहे.

\begin{figure}
  \begin{center}
    \begin{tikzpicture}[modal]
      \node[world] (w1) [label={[align=right]right:\mTrue{p}\\ \mFalse{q}}]{$w_1$};
      \node[world] (w2) [label={[align=right]right:\mTrue{p}\\ \mTrue{q}},
        above right=of w1]{$w_2$};
      \node[world] (w3) [label={[align=right]right:\mFalse{p}\\ \mFalse{q}},
        below right=of w1] {$w_3$};
      \draw[->] (w1) to (w2);
      \draw[->] (w1) to (w3);
    \end{tikzpicture}
  \end{center}
  \caption{साधे प्रतिमान.}
  \label{nml:syn:rel:fig:simple}
\end{figure}












\section{एखाद्या जगात सत्य असणे}\label{nml:syn:trw:sec}

प्रत्येक मोडल प्रतिमान त्यातील कोणती मोडल सूत्रे कोणत्या
जगांत सत्य मानायची हे ठरवते. ``प्रतिमान~$\mModel{M}$ मध्ये
सूत्र~$\varphi$ जग~$w$ येथे सत्य आहे'' हा संबंधी
चिन्हार्थमीमांसेतील मूलभूत संबंध आहे. हा संबंध विगमनाने
परिभाषित केला जातो; मोडल नसलेल्या परिकर्मींसाठी तो
सत्यताकोष्टकांनी दिलेल्या नेहमीच्या लक्षणाशी जुळतो.

\begin{defn}\label{nml:syn:trw:defn:mmodels}
  प्रतिमान~$\mModel M$ मध्ये \emph{सूत्र~$\varphi$ हे~$w$ येथे
  सत्य असणे}, म्हणजे $\mSat{M}{\varphi}[w]$, पुढीलप्रमाणे
  विगमनाने परिभाषित करतो:
  \begin{enumerate}
  \item \indcase{\varphi}{\lfalse}{$\mSat{M}{\lfalse}[w]$ कधीच नसते}.
  \item \indcase{\varphi}{\ltrue}{$\mSat{M}{\ltrue}[w]$ नेहमी असते}.
  \item $\mSat{M}{p}[w]$ तेव्हा आणि केवळ तेव्हाच, जेव्हा $w \in V(p)$.
  \item \indcase{\varphi}{\lnot \psi}{$\mSat{M}{\indfrm}[w]$
    तेव्हा आणि केवळ तेव्हाच, जेव्हा $\mSat/{M}{\psi}[w]$}.
  \item \indcase{\varphi}{(\psi \land \chi)}{$\mSat{M}{\indfrm}[w]$
    तेव्हा आणि केवळ तेव्हाच, जेव्हा $\mSat{M}{\psi}[w]$ आणि
    $\mSat{M}{\chi}[w]$}.
  \item \indcase{\varphi}{(\psi \lor \chi)}{$\mSat{M}{\indfrm}[w]$
    तेव्हा आणि केवळ तेव्हाच, जेव्हा $\mSat{M}{\psi}[w]$ किंवा
    $\mSat{M}{\chi}[w]$} (किंवा दोन्ही).
  \item \indcase{\varphi}{(\psi \lif \chi)}{$\mSat{M}{\indfrm}[w]$
    तेव्हा आणि केवळ तेव्हाच, जेव्हा $\mSat/{M}{\psi}[w]$ किंवा
    $\mSat{M}{\chi}[w]$}.
  \item \indcase{\varphi}{(\psi \liff \chi)}{$\mSat{M}{\indfrm}[w]$
    तेव्हा आणि केवळ तेव्हाच, जेव्हा $\mSat{M}{\psi}[w]$ आणि
    $\mSat{M}{\chi}[w]$ ही दोन्ही असतात, किंवा
    $\mSat{M}{\psi}[w]$ आणि $\mSat{M}{\chi}[w]$ यांपैकी कोणतेही नसते}.
  \item\label{nml:syn:trw:defn:sub:mmodels-box}
    \indcase{\varphi}{\Box \psi}{$\mSat{M}{\indfrm}[w]$ तेव्हा आणि
    केवळ तेव्हाच, जेव्हा $Rww'$ असलेल्या प्रत्येक $w' \in W$
    साठी $\mSat{M}{\psi}[w']$.}
  \item\label{nml:syn:trw:defn:sub:mmodels-diamond}
    \indcase{\varphi}{\Diamond \psi}{$\mSat{M}{\indfrm}[w]$ तेव्हा आणि
    केवळ तेव्हाच, जेव्हा $Rww'$ असलेल्या किमान एका $w' \in W$
    साठी $\mSat{M}{\psi}[w']$.}
  \end{enumerate}
\end{defn}

\ref{nml:syn:trw:defn:sub:mmodels-box} या खंडानुसार, $Rww'$ असलेले
एकही~$w'$ नसेल तेव्हा सूत्र~$\Box \psi$ हे~$w$ येथे
सत्य असते, हे लक्षात घ्या. अशा वेळी $\Box \psi$ हे~$w$ येथे
\emph{रिक्तपणे} सत्य असते. शिवाय $\psi$ सत्य नसतानाही
$\Box \psi$ ची~$w$ येथे पूर्ति होऊ शकते. $\psi$ हे~$w$ येथे
सत्य असल्यामुळेच $\Diamond \psi$ हेही~$w$ येथे सत्य असेल
असे नाही. मात्र $Rww$ असेल, उदाहरणार्थ $R$ परावर्ती असेल,
तर ते खरे ठरते. $Rww'$ असलेले $w'$ मुळीच नसेल, तर कोणत्याही
$\varphi$ साठी $\mSat/{M}{\Diamond \varphi}[w]$ असते.

\begin{prob}
  \ref{nml:syn:rel:fig:simple} मधील प्रतिमान पाहा. पुढीलपैकी
  कोणती विधाने खरी आहेत?
    \begin{enumerate}
    \item $\mSat{M}{q}[w_1]$;
    \item $\mSat{M}{\lnot q}[w_3]$;
    \item $\mSat{M}{p \lor q}[w_1]$;
    \item $\mSat{M}{\Box (p \lor q)}[w_1]$;
    \item $\mSat{M}{\Box q}[w_3]$;
    \item $\mSat{M}{\Box \bot}[w_3]$;
    \item $\mSat{M}{\Diamond q}[w_1]$;
    \item $\mSat{M}{\Box q}[w_1]$;
    \item $\mSat{M}{\lnot \Box\Box \lnot q}[w_1]$.
    \end{enumerate}
\end{prob}



\begin{prop}\label{nml:syn:trw:prop:dual}
  \begin{enumerate}
  \item $\mSat{M}{\Box \varphi}[w]$ तेव्हा आणि केवळ तेव्हाच, जेव्हा $\mSat{M}{\lnot\Diamond\lnot \varphi}[w]$.
  \item $\mSat{M}{\Diamond \varphi}[w]$ तेव्हा आणि केवळ तेव्हाच, जेव्हा $\mSat{M}{\lnot\Box\lnot \varphi}[w]$.
  \end{enumerate}
\end{prop}

\begin{proof}
  \begin{enumerate}
    \item $\mSat{M}{\lnot\Diamond\lnot \varphi}[w]$ तेव्हा आणि केवळ
      तेव्हाच, जेव्हा $\mSat/{M}{\Diamond\lnot \varphi}[w]$;
      हे $\mSat{M}{}[w]$ च्या व्याख्येवरून मिळते.
      $\mSat{M}{\Diamond\lnot \varphi}[w]$ तेव्हा आणि केवळ तेव्हाच,
      जेव्हा $Rww'$ असलेल्या एखाद्या $w'$ साठी
      $\mSat{M}{\lnot \varphi}[w']$. म्हणून
      $\mSat/{M}{\Diamond\lnot \varphi}[w]$ तेव्हा आणि केवळ तेव्हाच,
      जेव्हा $Rww'$ असलेल्या प्रत्येक $w'$ साठी
      $\mSat/{M}{\lnot \varphi}[w']$. तसेच
      $\mSat/{M}{\lnot \varphi}[w']$ तेव्हा आणि केवळ तेव्हाच,
      जेव्हा $\mSat{M}{\varphi}[w']$. मिळून,
      $\mSat{M}{\lnot\Diamond\lnot \varphi}[w]$ तेव्हा आणि केवळ
      तेव्हाच, जेव्हा $Rww'$ असलेल्या प्रत्येक $w'$ साठी
      $\mSat{M}{\varphi}[w']$. $\mSat{M}{}[w]$ च्या व्याख्येवरून
      ही अट तेव्हा आणि केवळ तेव्हाच असते, जेव्हा
      $\mSat{M}{\Box \varphi}[w]$.
    \item $\mSat{M}{\lnot\Box\lnot \varphi}[w]$
      तेव्हा आणि केवळ तेव्हाच, जेव्हा
      $\mSat/{M}{\Box\lnot \varphi}[w]$.
      $\mSat{M}{\Box\lnot \varphi}[w]$ तेव्हा आणि केवळ तेव्हाच,
      जेव्हा $Rww'$ असलेल्या प्रत्येक $w'$ साठी
      $\mSat{M}{\lnot \varphi}[w']$. म्हणून
      $\mSat/{M}{\Box\lnot \varphi}[w]$ तेव्हा आणि केवळ तेव्हाच,
      जेव्हा $Rww'$ असलेल्या एखाद्या $w'$ साठी
      $\mSat/{M}{\lnot \varphi}[w']$. तसेच
      $\mSat/{M}{\lnot \varphi}[w']$ तेव्हा आणि केवळ तेव्हाच,
      जेव्हा $\mSat{M}{\varphi}[w']$. मिळून,
      $\mSat{M}{\lnot\Box\lnot \varphi}[w]$ तेव्हा आणि केवळ
      तेव्हाच, जेव्हा $Rww'$ असलेल्या एखाद्या~$w'$ साठी
      $\mSat{M}{\varphi}[w']$. पुन्हा $\mSat{M}{}[w]$ च्या
      व्याख्येवरून ही अट तेव्हा आणि केवळ तेव्हाच असते,
      जेव्हा $\mSat{M}{\Diamond \varphi}[w]$.
  \end{enumerate}
\end{proof}




\begin{prob}
  $\mModel{M} = \tuple{W, R, V}$ हे प्रतिमान असू द्या आणि
  $w_1, w_2 \in W$ साठी पुढील अटी पूर्ण होतात असे समजा:
  \begin{enumerate}
  \item प्रत्येक विधानीय चल~$p$ साठी $w_1 \in V(p)$
    तेव्हा आणि केवळ तेव्हाच, जेव्हा $w_2 \in V(p)$; आणि
  \item प्रत्येक $w \in W$ साठी: $Rw_1w$ तेव्हा आणि केवळ
    तेव्हाच, जेव्हा $Rw_2w$.
  \end{enumerate}
  सूत्रे वर विगमन करून दाखवा की सर्व सूत्रे~$\varphi$
  साठी $\mSat{M}{\varphi}[w_1]$ तेव्हा आणि केवळ तेव्हाच,
  जेव्हा $\mSat{M}{\varphi}[w_2]$.
\end{prob}

\begin{prob}
  $\mModel{M} = \tuple{W, R, V}$ असू द्या. दाखवा की
  $\mSat{M}{\lnot\Diamond \varphi}[w]$ तेव्हा आणि केवळ तेव्हाच,
  जेव्हा $\mSat{M}{\Box\lnot\varphi}[w]$.
\end{prob}










\section{प्रतिमानात सत्य असणे}\label{nml:syn:tru:sec}

कधीकधी दिलेल्या प्रतिमानातील प्रत्येक जगात कोणती सूत्रे सत्य
आहेत, हे जाणून घ्यायचे असते. यासाठी एक संकेतन देऊ.

\begin{defn}
  सूत्र~$\varphi$ हे $M = \tuple{W, R,
    V}$ या \emph{प्रतिमानात सत्य आहे}, म्हणजे $\mSat{M}{\varphi}$,
  तेव्हा आणि केवळ तेव्हाच, जेव्हा प्रत्येक $w \in W$ साठी
  $\mSat{M}{\varphi}[w]$.
\end{defn}

\begin{prop}\label{nml:syn:tru:prop:truthfacts}
  \begin{enumerate}
  \item $\mSat{M}{\varphi}$ असेल, तर $\mSat/{M}{\lnot \varphi}$;
    पण याचा \emph{व्यत्यास नाही}.
  \item $\mSat{M}{\varphi \lif \psi}$ असेल, तर $\mSat{M}{\varphi}$
    फक्त $\mSat{M}{\psi}$ असेल तरच; पण याचा \emph{व्यत्यास नाही}.
  \end{enumerate}
\end{prop}

\begin{proof}
  \begin{enumerate}
  \item $\mSat{M}{\varphi}$ असेल, तर $\varphi$ हे $W$ मधील सर्व
    जगांत सत्य आहे. $W \neq \emptyset$ असल्यामुळे
    $\mSat{M}{\lnot\varphi}$ असू शकत नाही; तसे असल्यास
    एखाद्या जगात $\varphi$ सत्यही आणि असत्यही असावे लागेल.

    दुसरीकडे, $\mSat/{M}{\lnot \varphi}$ असेल, तर एखाद्या
    $w \in W$ येथे $\varphi$ सत्य आहे. यावरून \emph{प्रत्येक}
    $w \in W$ साठी $\mSat{M}{\varphi}[w]$ मिळत नाही. उदाहरणार्थ,
    \ref{nml:syn:rel:fig:simple} मधील प्रतिमानात $\mSat/{M}{\lnot p}$
    आणि $\mSat/{M}{p}$ ही दोन्ही आहेत.
  \item $\mSat{M}{\varphi \lif \psi}$ आणि $\mSat{M}{\varphi}$ असे समजा.
    $\mSat{M}{\psi}$ दाखवण्यासाठी कोणतेही $w \in W$ घ्या.
    मग $\mSat{M}{\varphi \lif \psi}[w]$ आणि $\mSat{M}{\varphi}[w]$;
    म्हणून $\mSat{M}{\psi}[w]$. $w$ कोणतेही असल्यामुळे
    $\mSat{M}{\psi}$.

    व्यत्यास असत्य ठरवण्यासाठी असे प्रतिमान $\mModel{M}$
    शोधावे लागेल की $\mSat{M}{\varphi}$ हे फक्त $\mSat{M}{\psi}$
    असेल तरच, पण $\mSat/{M}{\varphi \lif \psi}$. पुन्हा
    \ref{nml:syn:rel:fig:simple} मधील प्रतिमान पाहा: $\mSat/{M}{p}$,
    आणि म्हणून $\mSat{M}{p}$ हे फक्त $\mSat{M}{q}$ असेल
    तरच, हे रिक्तपणे खरे आहे. मात्र $w_1$ येथे $p$ सत्य
    आणि $q$ असत्य असल्यामुळे $\mSat/{M}{p \lif
      q}$.
  \end{enumerate}
\end{proof}

\begin{prob}
  $p_1$, $p_2$, $p_3$ हीच विधानीय चले असलेल्या
  भाषेसाठी पुढील प्रतिमान $\mModel{M}$ पाहा:
  \begin{center}
    \begin{tikzpicture}[modal]
      \node[world] (w1) [label={[align=right]left:\mTrue{p_1}\\\mFalse{p_2}\\\mFalse{p_3}}]
            {$w_1$} ;
      \node[world] (w2) [label={[align=right]right:\mTrue{p_1}\\\mTrue{p_2}\\\mFalse{p_3}},
        below right=of w1]  {$w_2$};
      \node[world] (w3) [label={[align=right]right:\mTrue{p_1}\\\mTrue{p_2}\\\mTrue{p_3}},
        above right=of w2] {$w_3$};
      \draw[reflexive above] (w3) to (w3);
      \draw[->] (w1) to (w2);
      \draw[->] (w2) to (w3);
      \draw[->] (w1) to (w3);
    \end{tikzpicture}
  \end{center}
  पुढील सूत्रे आणि रूपबंध $\mModel{M}$ या प्रतिमानात,
  म्हणजे $\mModel{M}$ मधील प्रत्येक जगात, सत्य आहेत का? स्पष्ट करा.
  \begin{enumerate}
  \item $p\lif \Diamond p$ ($p$ अणुरूप असेल तेव्हा);
  \item $\varphi\lif \Diamond \varphi$ ($\varphi$ कोणतेही असेल तेव्हा);
  \item $\Box p \lif p$ ($p$ अणुरूप असेल तेव्हा);
  \item $\lnot p \lif \Diamond \Box p$ ($p$ अणुरूप असेल तेव्हा);
  \item $\Diamond \Box \varphi$ ($\varphi$ कोणतेही असेल तेव्हा);
  \item $\Box \Diamond p$ ($p$ अणुरूप असेल तेव्हा).
  \end{enumerate}
\end{prob}












\section{वैधता}\label{nml:syn:val:sec}

\begin{explain}
  सर्व प्रतिमानांत, म्हणजे प्रत्येक प्रतिमानातील प्रत्येक जगात, सत्य
  असणारी सूत्रे विशेष महत्त्वाची आहेत. संभाव्य जगांवरील
  एखाद्या प्राप्यता संबंधापासून अर्थलक्षण मिळते, या अर्थाने
  $\Box$ आणि $\Diamond$ यांचे अर्थलक्षण ``सामान्य'' असेल,
  तर ते कोणतेही असो, सत्य असणारी मोडल विधाने ही सूत्रे व्यक्त
  करतात. अशा सूत्रे ना आपण \emph{वैध} म्हणतो. उदाहरणार्थ,
  $\Box(p \land q) \lif \Box p$ हे वैध आहे. मात्र $\Box$
  च्या अनिवार्यतेशी संबंधित अर्थावरून वैध वाटणारी काही
  सूत्रे, उदाहरणार्थ $\Box p \lif p$, वैध नसतात.
  संबंधी प्रतिमानांचे एक महत्त्व असे की, वेगवेगळ्या प्रकारच्या
  प्राप्यता संबंधांनी $\Box$ आणि $\Diamond$ यांचे वेगवेगळे
  अर्थलक्षण व्यक्त करता येते. त्यामुळे वैधता केवळ \emph{सर्व}
  प्रतिमानांच्या संदर्भात नव्हे, तर \emph{विशिष्ट प्रकारच्या}
  सर्व प्रतिमानांच्या संदर्भातही परिभाषित करावी. उदाहरणार्थ,
  ज्या प्रतिमानांत प्रत्येक जग स्वतःलाच प्राप्य आहे, म्हणजे
  $R$ परावर्ती आहे, त्या सर्व प्रतिमानांत $\Box p \lif p$
  सत्य आहे, हे पुढे दिसेल. प्रतिमानांच्या वर्गांच्या संदर्भात
  वैधता परिभाषित केल्यास हे थोडक्यात सांगता येते:
  $\Box p \lif p$ हे परावर्ती प्रतिमानांच्या वर्गात वैध आहे.
\end{explain}

\begin{defn}
  सूत्र~$\varphi$ हे प्रतिमानांच्या $\mClass{C}$ या वर्गात
  \emph{वैध} आहे, जर ते $\mClass{C}$ मधील प्रत्येक प्रतिमानात
  सत्य असेल (म्हणजे $\mClass{C}$ मधील प्रत्येक प्रतिमानातील
  प्रत्येक जगात सत्य असेल). $\varphi$ हे $\mClass{C}$ मध्ये वैध
  असेल, तर $\mClass{C} \Entails \varphi$ असे लिहितो; आणि $\varphi$
  हे \emph{सर्व} प्रतिमानांच्या वर्गात वैध असेल, तर
  $\Entails \varphi$ असे लिहितो.
\end{defn}

\begin{prop}\label{nml:syn:val:prop:subset-class}
  $\varphi$ हे $\mClass{C}$ मध्ये वैध असेल, तर प्रत्येक
  $\mClass{C}' \subseteq \mClass{C}$ या वर्गातही ते वैध असते.
\end{prop}

\begin{prop}\label{nml:syn:val:prop:Nec-rule}
  $\varphi$ वैध असेल, तर $\Box\varphi$ देखील वैध असते.
\end{prop}

\begin{proof}
  $\Entails \varphi$ असे समजा. $\Entails \Box\varphi$ दाखवण्यासाठी
  कोणतेही प्रतिमान $\mModel{M} = \tuple{W, R, V}$ आणि
  $w \in W$ घ्या. $Rww'$ असेल, तर $\varphi$ वैध असल्यामुळे
  $\mSat{M}{\varphi}[w']$; म्हणून $\mSat{M}{\Box\varphi}[w]$ देखील.
  $\mModel{M}$ आणि $w$ कोणतेही असल्यामुळे
  $\Entails \Box\varphi$.
\end{proof}

\begin{prob}
  पुढील सूत्रे वैध आहेत हे दाखवा:
  \begin{enumerate}
  \item $\Box p \lif \Box (q \lif p)$;
  \item $\Box \lnot \lfalse$;
  \item $\Box p \lif (\Box q \lif \Box p)$.
  \end{enumerate}
\end{prob}

\begin{prob}
  $W = \{w\}$ असलेल्या $\mModel{M} = \tuple{W, R, V}$
  प्रतिमानांच्या $\mClass{C}$ या वर्गात $\varphi \lif \Box\varphi$
  वैध आहे, हे दाखवा. तसेच $R = \emptyset$ असलेल्या
  $\mModel{M} = \tuple{W, R, V}$ प्रतिमानांच्या वर्गात
  $\psi \lif \Box \varphi$ आणि $\Diamond \varphi \lif \psi$ वैध
  आहेत, हे दाखवा.
\end{prob}












\section{सर्वतः सत्य सूत्रांचे आदेशन-दृष्टांत}\label{nml:syn:tau:sec}

\begin{explain}
मोडल परिकर्मी नसलेले सूत्र प्रत्येक सत्यतामूल्य-नियुक्तीखाली
सत्य असेल, तर ते सर्वतः सत्य सूत्र असते. असे प्रत्येक सूत्र
प्रत्येक प्रतिमानातील प्रत्येक जगात सत्य असते, हे स्पष्ट आहे.
पण $\Box$ आणि~$\Diamond$ असलेल्या सूत्रे साठी
सर्वतः सत्य असण्याची संकल्पना परिभाषित केलेली नाही.
उदाहरणार्थ, विमध्य सिद्धांताचे दृष्टांत सूत्र असलेले
$\Box p \lor \lnot \Box p$ वैध आहे का? येथे
\emph{सर्वतः सत्य सूत्राचा आदेशन-दृष्टांत} ही
संकल्पना उपयोगी पडते: ते मोडल नसलेल्या सर्वतः सत्य
सूत्रात आदेशन करून मिळणारे सूत्र असते. असे
प्रत्येक आदेशन-दृष्टांत सूत्र वैध असणे अपेक्षित आहे;
तरीही ते सिद्ध करावे लागते.
\end{explain}

\begin{defn}
  मोडल सूत्र~$\psi$ हे \emph{सर्वतः सत्य सूत्राचा
  आदेशन-दृष्टांत} तेव्हा आणि केवळ तेव्हाच असते,
  जेव्हा $p_1$, \dots,~$p_n$ ही विधानीय चले असलेले मोडल नसलेले सर्वतः सत्य सूत्र~$\varphi$
  आणि सूत्रे $\theta_1$, \dots,~$\theta_n$ अशी अस्तित्वात
  असतात की $\psi \ident \SSubst{\varphi}{\subst{\theta_1}{p_1},
    \dots, \subst{\theta_n}{p_n}}$.
\end{defn}

\begin{lem}\label{nml:syn:tau:lem:valid-taut}
  $\varphi$ हे मोडल नसलेले सूत्र असून त्यातील
  विधानीय चले $p_1$, \dots, $p_n$ आहेत,
  आणि $\theta_1$, \dots, $\theta_n$ ही मोडल सूत्रे आहेत असे
  समजा. मग कोणतीही नियुक्ती $\pAssign{v}$, कोणतेही प्रतिमान
  $\mModel{M} = \tuple{W, R, V}$, आणि असे कोणतेही $w \in W$
  घ्या की $\pAssign{v}(p_i) = \True$ तेव्हा आणि केवळ तेव्हाच,
  जेव्हा $\mSat{M}{\theta_i}[w]$. अशा वेळी $\pSat{v}{\varphi}$
  तेव्हा आणि केवळ तेव्हाच, जेव्हा
  $\mSat{M}{\SSubst{\varphi}{\subst{\theta_1}{p_1}, \dots,
      \subst{\theta_n}{p_n}}}[w]$.
\end{lem}

\begin{proof}
  $\varphi$ वर विगमन करू.
  \begin{enumerate}
  \item \indcase{\varphi}{\lfalse}{$\pSat/{v}{\lfalse}$
      आणि $\mSat/{M}{\lfalse}[w]$ ही दोन्ही असतात}.
  \item \indcase{\varphi}{\ltrue}{$\pSat{v}{\ltrue}$
      आणि $\mSat{M}{\ltrue}[w]$ ही दोन्ही असतात}.
  \item \indcase{\varphi}{p_i}{
    \begin{align*}
      \pSat{v}{p_i} \Leftrightarrow {} & \pAssign{v}(p_i) = \True \\
      & \qquad \text{$\pSat{v}{p_i}$ च्या व्याख्येवरून}\\
      \Leftrightarrow {} & \mSat{M}{\theta_i}[w] \\
      &\qquad \text{गृहीतकावरून}\\
      \Leftrightarrow {} & \mSat{M}{\SSubst{p_i}{\subst{\theta_1}{p_1}, \dots,
          \subst{\theta_n}{p_n}}}[w]\\
      &\qquad \text{$\SSubst{p_i}{\subst{\theta_1}{p_1}, \dots,
          \subst{\theta_n}{p_n}} \ident \theta_i$ असल्यामुळे}.
      \end{align*}}
  \item \indcase{\varphi}{\lnot \psi}{
      \begin{align*}
        \pSat{v}{\lnot \psi} \Leftrightarrow {} & \pSat/{v}{\psi}\\
        &\qquad \text{$\pSat{v}{}$ च्या व्याख्येवरून};\\
        \Leftrightarrow {} & \mSat/{M}{\SSubst{\psi}{\subst{\theta_1}{p_1}, \dots,
          \subst{\theta_n}{p_n}}}[w]\\
        &\qquad \text{विगमन गृहीतकावरून}\\
        \Leftrightarrow {} &
        \mSat{M}{\SSubst{\lnot \psi}{\subst{\theta_1}{p_1}, \dots,
          \subst{\theta_n}{p_n}}}[w]\\
        &\qquad \text{मोडल पूर्तिच्या $\mSat{M}{}[w]$ व्याख्येवरून}.
      \end{align*}}
  \item \indcase{\varphi}{(\psi \land \chi)}{
      \begin{align*}
        \pSat{v}{\psi \land \chi} \Leftrightarrow {} &
        \pSat{v}{\psi} \text{ आणि } \pSat{v}{\chi}\\
        &\qquad \text{$\pSat{v}{}$ च्या व्याख्येवरून}\\
        \Leftrightarrow {} &
        \mSat{M}{\SSubst{\psi}{\subst{\theta_1}{p_1}, \dots,
            \subst{\theta_n}{p_n}}}[w] \text{ आणि } \\
        & \mSat{M}{\SSubst{\chi}{\subst{\theta_1}{p_1}, \dots,
          \subst{\theta_n}{p_n}}}[w]\\
        &\qquad \text{विगमन गृहीतकावरून}\\
        \Leftrightarrow{} &
        \mSat{M}{\SSubst{(\psi \land \chi)}{\subst{\theta_1}{p_1}, \dots,
          \subst{\theta_n}{p_n}}}[w]\\
        &\qquad \text{$\mSat{M}{}[w]$ च्या व्याख्येवरून}.
      \end{align*}}
  \item \indcase{\varphi}{(\psi \lor \chi)}{
      \begin{align*}
        \pSat{v}{\psi \lor \chi} \Leftrightarrow{} &
        \pSat{v}{\psi} \text{ किंवा } \pSat{v}{\chi}\\
        &\qquad \text{$\pSat{v}{}$ च्या व्याख्येवरून};\\
        \Leftrightarrow{} & \mSat{M}{\SSubst{\psi}{\subst{\theta_1}{p_1}, \dots,
            \subst{\theta_n}{p_n}}}[w] \text{ किंवा }\\
        & \mSat{M}{\SSubst{\chi}{\subst{\theta_1}{p_1}, \dots,
          \subst{\theta_n}{p_n}}}[w]\\
        &\qquad \text{विगमन गृहीतकावरून}\\
        \Leftrightarrow{} &
        \mSat{M}{\SSubst{(\psi \lor \chi)}{\subst{\theta_1}{p_1}, \dots,
          \subst{\theta_n}{p_n}}}[w]\\
        &\qquad \text{$\mSat{M}{}[w]$ च्या व्याख्येवरून}.
      \end{align*}}
  \item \indcase{\varphi}{(\psi \lif \chi)}{
      \begin{align*}
        \pSat{v}{\psi \lif \chi} \Leftrightarrow{} &
        \pSat/{v}{\psi} \text{ किंवा } \pSat{v}{\chi}\\
        &\qquad \text{$\pSat{v}{}$ च्या व्याख्येवरून}\\
        \Leftrightarrow{} &
        \mSat/{M}{\SSubst{\psi}{\subst{\theta_1}{p_1}, \dots,
            \subst{\theta_n}{p_n}}}[w] \text{ किंवा }\\
        & \mSat{M}{\SSubst{\chi}{\subst{\theta_1}{p_1}, \dots,
          \subst{\theta_n}{p_n}}}[w]\\
        &\qquad \text{विगमन गृहीतकावरून}\\
        \Leftrightarrow{} &
        \mSat{M}{\SSubst{(\psi \lif \chi)}{\subst{\theta_1}{p_1}, \dots,
          \subst{\theta_n}{p_n}}}[w]\\
        &\qquad \text{$\mSat{M}{}[w]$ च्या व्याख्येवरून}.
      \end{align*}}
  \item \indcase{\varphi}{(\psi \liff \chi)}{
      \begin{align*}
        \pSat{v}{\psi \liff \chi} \Leftrightarrow {} &
        \text{एकतर } \pSat{v}{\psi} \text{ आणि } \pSat{v}{\chi}\\
        & \text{किंवा } \pSat/{v}{\psi} \text{ आणि } \pSat/{v}{\chi}\\
        &\qquad \text{$\pSat{v}{}$ च्या व्याख्येवरून}\\
        \Leftrightarrow {} &
        \text{एकतर } \mSat{M}{\SSubst{\psi}{\subst{\theta_1}{p_1}, \dots,
            \subst{\theta_n}{p_n}}}[w] \text{ आणि }\\
        & \qquad \mSat{M}{\SSubst{\chi}{\subst{\theta_1}{p_1}, \dots,
          \subst{\theta_n}{p_n}}}[w]\\
        &  \text{ किंवा } \mSat/{M}{\SSubst{\psi}{\subst{\theta_1}{p_1}, \dots,
            \subst{\theta_n}{p_n}}}[w] \text{ आणि }\\
        & \qquad \mSat/{M}{\SSubst{\chi}{\subst{\theta_1}{p_1}, \dots,
          \subst{\theta_n}{p_n}}}[w]\\
        &\qquad\text{विगमन गृहीतकावरून}\\
        \Leftrightarrow {} &
        \mSat{M}{\SSubst{(\psi \liff \chi)}{\subst{\theta_1}{p_1}, \dots,
          \subst{\theta_n}{p_n}}}[w]\\
        &\qquad \text{$\mSat{M}{}[w]$ च्या व्याख्येवरून}.
      \end{align*}}
  \end{enumerate}
  \end{proof}

\begin{prop}\label{nml:syn:tau:prop:valid-taut}
  सर्वतः सत्य सूत्रांचे सर्व आदेशन-दृष्टांत वैध असतात.
\end{prop}

\begin{proof}
  प्रतिपरिवर्तनाने सिद्ध करू. काही प्रतिमान $\mModel{M}$ आणि जग~$w$
  यांसाठी पुढील अट असलेले $\varphi$ आहे असे समजा:
  $\mSat/{M}{\SSubst{\varphi}{\subst{\theta_1}{p_1}, \dots,
  \subst{\theta_n}{p_n}}}[w]$.
  $\pAssign{v}(p_i) = \True$ तेव्हा आणि केवळ तेव्हाच,
  जेव्हा $\mSat{M}{\theta_i}[w]$, अशी नियुक्ती $\pAssign{v}$
  परिभाषित करा ($q \notin \{p_1, \dots, p_n \}$ या
  चलांना $\pAssign{v}$ कोणतीही मूल्ये देते). मग
  \ref{nml:syn:tau:lem:valid-taut} वरून $\pSat/{v}{\varphi}$; म्हणून
  $\varphi$ सर्वतः सत्य सूत्र नाही.
\end{proof}













\section{रूपबंध आणि वैधता}\label{nml:syn:sch:sec}

\begin{defn}
  \emph{रूपबंध} म्हणजे एखाद्या मोडल सूत्र~$\chi$ चे
  नेमके सर्व आदेशन-दृष्टांत मिळून बनलेला सूत्रांचा संच; म्हणजे
  \[  \Setabs{\psi}{\lexists[\theta_1], \dots, \lexists[\theta_n] \left(\psi =
  \SSubst{\chi}{\subst{\theta_1}{p_1}, \dots, \subst{\theta_n}{p_n}} \right) }.
  \]
  सूत्र~$\chi$ याला रूपबंधाचे \emph{वैशिष्ट्यदर्शक} सूत्र
  म्हणतात. विधानीय चलांची नावे बदलण्याचा फरक सोडल्यास ते एकमेव असते.
  सूत्र~$\varphi$ हे त्या संचाचा घटक असेल, तर तो
  रूपबंधाचा \emph{दृष्टांत} असतो.
\end{defn}

रूपबंध दर्शवण्यासाठी $\chi$ मधील विधानीय चलांच्या जागी
`$\varphi$', `$\psi$',~\dots ही चिन्हे बसवून मिळणारी
अधिभाषेतील अभिव्यक्ती वापरणे सोयीचे असते. तार्किक स्थिरांक बदलत नाहीत. उदाहरणार्थ,
`$\varphi$', `$\varphi \lif \Box\varphi$', `$\varphi \lif (\psi \lif \varphi)$'
ही रूपबंध दर्शवतात. त्यांची वैशिष्ट्यदर्शक सूत्रे
अनुक्रमे $p$, $p \lif \Box p$, $p \lif (q
\lif p)$ आहेत. `$\varphi$' हा रूपबंध \emph{सर्व}
सूत्रांचा संच दर्शवतो.

\begin{defn}
  रूपबंधाचे सर्व दृष्टांत एखाद्या प्रतिमानात सत्य असतील तेव्हा
  आणि केवळ तेव्हाच तो रूपबंध त्या प्रतिमानात \emph{सत्य}
  असतो; आणि प्रत्येक प्रतिमानात सत्य असेल तेव्हा आणि केवळ
  तेव्हाच तो \emph{वैध} असतो.
\end{defn}

\begin{prop}\label{nml:syn:sch:prop:Kvalid}
  पुढील \Ax{K} हा रूपबंध वैध आहे:
  \begin{equation*}
  \Box(\varphi \lif \psi) \lif (\Box \varphi \lif \Box \psi). \tag{\Ax{K}}
  \end{equation*}
\end{prop}

\begin{proof}
  या रूपबंधाचा प्रत्येक दृष्टांत प्रत्येक प्रतिमानातील प्रत्येक
  जगात सत्य आहे, हे दाखवायचे आहे. म्हणून कोणतेही
  $\mModel{M} = \tuple{W,R,V}$ आणि $w \in W$ घ्या.
  एखादे अभिव्यंजन एखाद्या जगात सत्य असल्याचे दाखवताना,
  त्याचे पूर्वांग सत्य मानून उत्तरांगही सत्य आहे हे दाखवतो.
  येथे $\mSat{M}{\Box(\varphi \lif \psi)}[w]$ आणि
  $\mSat{M}{\Box \varphi}[w]$ असू द्या. आता
  $\mSat{M}{\Box \psi}[w]$ दाखवायचे आहे. म्हणून $Rww'$
  असलेले कोणतेही $w'$ घ्या. पहिल्या गृहीतकावरून
  $\mSat{M}{\varphi \lif \psi}[w']$ आणि दुसऱ्यावरून
  $\mSat{M}{\varphi}[w']$. त्यामुळे $\mSat{M}{\psi}[w']$.
  $w'$ कोणतेही असल्यामुळे $\mSat{M}{\Box \psi}[w]$.
\end{proof}

\begin{prop}\label{nml:syn:sch:prop:Dual-valid}
  पुढील \Dual{} हा रूपबंध वैध आहे:
  \begin{equation*}
  \Diamond \varphi \liff \lnot\Box\lnot \varphi. \tag{\Dual}
  \end{equation*}
\end{prop}

\begin{proof}
  सराव.
\end{proof}

\begin{prob}
  \ref{nml:syn:sch:prop:Dual-valid} सिद्ध करा.
\end{prob}

\begin{prop}\label{nml:syn:sch:prop:soundMP}
  प्रतिमानातील एखाद्या जगात $\varphi$ आणि $\varphi \lif \psi$ सत्य
  असतील, तर $\psi$ देखील सत्य असते. म्हणून वैध सूत्रे
  विध्यात्मक अनुमानाखाली बंद असतात.
\end{prop}

\begin{prop}\label{nml:syn:sch:prop:valid-instances}
  सूत्र~$\varphi$ तेव्हा आणि केवळ तेव्हाच वैध असते,
  जेव्हा त्याचे सर्व आदेशन-दृष्टांत वैध असतात. दुसऱ्या शब्दांत,
  रूपबंध तेव्हा आणि केवळ तेव्हाच वैध असतो, जेव्हा त्याचे
  वैशिष्ट्यदर्शक सूत्र वैध असते.
\end{prop}

\begin{proof}
  ``जर'' ही दिशा स्पष्ट आहे, कारण $\varphi$ स्वतःचाच
  आदेशन-दृष्टांत आहे.

  ``फक्त तरच'' ही दिशा सिद्ध करण्यासाठी पुढील दावा दाखवू.
  $\mModel{M} = \tuple{W, R, V}$ हे मोडल प्रतिमान असू द्या,
  आणि $\psi \ident
  \SSubst{\varphi}{\subst{\theta_1}{p_1},\dots,\subst{\theta_n}{p_n}}$ हा
  $\varphi$ चा आदेशन-दृष्टांत असू द्या. $1 \leq i \leq n$ साठी
  $V'(p_i) = \Setabs{w \in W}{\mSat{M}{\theta_i}[w]}$ ठेवा;
  इतर प्रत्येक विधानीय चल~$q \notin \{p_1,\dots,p_n\}$
  साठी $V'(q)=V(q)$ ठेवा. या सर्व चलांवरील नियुक्तीने
  $\mModel{M'} = \tuple{W, R,V'}$ परिभाषित करा.
  मग कोणत्याही~$w \in W$ साठी
  $\mSat{M}{\psi}[w]$ तेव्हा आणि केवळ तेव्हाच, जेव्हा
  $\mSat{M'}{\varphi}[w]$. (याची सिद्धता सरावासाठी ठेवली आहे.)
  आता $\varphi$ वैध आहे, पण $\varphi$ चा एखादा आदेशन-दृष्टांत
  $\psi$ वैध नाही, असे समजा. मग काही
  $\mModel{M} = \tuple{W, R, V}$ आणि काही $w \in W$
  साठी $\mSat/{M}{\psi}[w]$. पण दाव्यावरून
  $\mSat/{M'}{\varphi}[w]$; म्हणजे $\varphi$ वैध नाही. हा विरोधाभास.
\end{proof}

\begin{prob}
  \ref{nml:syn:sch:prop:valid-instances} च्या सिद्धतेतील
  ``फक्त तरच'' भागातला दावा सिद्ध करा. (सूचना:
  $\varphi$ वर विगमन करा.)
\end{prob}

मात्र एखादा रूपबंध प्रतिमानात सत्य असणे आणि त्याचे
वैशिष्ट्यदर्शक सूत्र त्या प्रतिमानात सत्य असणे या सममूल्य
अटी नाहीत, हे लक्षात घ्या. ``फक्त तरच'' ही दिशा खरी आहे:
$\varphi$ चा प्रत्येक दृष्टांत~$\mModel{M}$ मध्ये सत्य असेल,
तर $\varphi$~स्वतःही~$\mModel{M}$ मध्ये सत्य असते. पण
$\varphi$ हे~$\mModel{M}$ मध्ये सत्य असूनही $\varphi$ चा एखादा
दृष्टांत~$\mModel{M}$ मधील काही जगात असत्य असू शकतो.
अगदी साध्या प्रतिदृष्टांतासाठी~$p$ घेऊन एकच जग~$w$ असलेले आणि
$V(p) = \{w\}$ असलेले प्रतिमान घ्या. तेथे~$p$ हे~$w$
येथे सत्य आहे. पण $\lfalse$ हा~$p$ चा दृष्टांत असून
तो~$w$ येथे सत्य नाही.

\begin{prob}
  पुढीलपैकी कोणतेही सूत्रे वैध नाही, हे दाखवा:
  \begin{enumerate}
    \item[\Ax{D}:] \quad $\Box p \lif \Diamond p$;
    \item[\Ax{T}:] \quad $\Box p \lif p$;
    \item[\Ax{B}:] \quad $p \lif \Box\Diamond p$;
    \item[\Ax{4}:] \quad $\Box p \lif \Box \Box p$;
    \item[\Ax{5}:] \quad $\Diamond p \lif \Box \Diamond
      p$.
  \end{enumerate}
\end{prob}

\begin{table}[t]
    \centering
    \begin{tabular}{| l || l |}
      \hline
      {\emph{वैध रूपबंध}} & {\emph{अवैध रूपबंध}} \\
      \hline\hline
      $\Box(\varphi \lif \psi) \lif (\Diamond \varphi \lif \Diamond \psi)$
      & $\Box (\varphi \lor \psi) \lif (\Box \varphi \lor \Box \psi)$ \\
      $\Diamond (\varphi \lif \psi) \lif (\Box \varphi \lif \Diamond
      \psi)$
      & $(\Diamond \varphi \land \Diamond \psi) \lif \Diamond (\varphi
      \land \psi)$\\
      $\Box (\varphi \land \psi) \liff (\Box \varphi \land \Box \psi)$
      & $\varphi \lif \Box \varphi$ \\
      $\Box \varphi \lif \Box (\psi \lif \varphi)$
      & $\Box \Diamond \varphi \lif \psi$ \\
      $\lnot \Diamond \varphi \lif \Box (\varphi \lif \psi)$
      & $\Box \Box \varphi \lif \Box \varphi$ \\
      $\Diamond (\varphi \lor \psi) \liff (\Diamond \varphi \lor
      \Diamond \psi)$
      & $\Box \Diamond \varphi \lif \Diamond \Box \varphi$. \\
      \hline
    \end{tabular}
    \caption{वैध आणि अवैध रूपबंध.}
    \label{nml:syn:sch:tab:valid-invalidSchemas}
\end{table}

\begin{prob}
  \ref{nml:syn:sch:tab:valid-invalidSchemas} च्या
  पहिल्या स्तंभातील रूपबंध वैध आहेत, आणि दुसऱ्या
  स्तंभातील रूपबंध वैध नाहीत, हे सिद्ध करा.
\end{prob}

\begin{prob}
  पुढील रूपबंध वैध आहेत की अवैध, ते ठरवा:
  \begin{enumerate}
  \item $(\Diamond \varphi \lif \Box \psi) \lif (\Box \varphi \lif \Box
    \psi)$;
  \item $\Diamond(\varphi \lif \psi) \lor \Box(\psi \lif \varphi)$.
  \end{enumerate}
\end{prob}

\begin{prob}
  पुढील प्रत्येक रूपबंधासाठी असे प्रतिमान $\mModel{M}$ शोधा की
  त्या सूत्राचा प्रत्येक दृष्टांत $\mModel{M}$ मध्ये सत्य आहे:
  \begin{enumerate}
  \item $p \lif \Diamond\Diamond p$;
  \item $\Diamond p \lif \Box p$.
  \end{enumerate}
\end{prob}












\section{तार्किक निष्पन्नता}\label{nml:syn:ent:sec}

\begin{explain}
  एखाद्या जगात सत्य असण्याच्या व्याख्येने सूत्रे मधील
  तार्किक निष्पन्नतेचा संबंध परिभाषित करता येतो.
  सूत्र~$\psi$ पासून~$\varphi$ तार्किकरीत्या निष्पन्न
  होते तेव्हा आणि केवळ तेव्हाच, जेव्हा $\psi$ सत्य असेल
  तेव्हा $\varphi$ देखील सत्य असते. येथे ``तेव्हा'' म्हणताना
  कोणतेही प्रतिमान आणि त्यातील कोणतेही जग अभिप्रेत आहे.
\end{explain}

\begin{defn}
  $\Gamma$ हा सूत्रांचा संच आणि $\varphi$ हे सूत्र
  असेल, तर $\Gamma$ पासून~$\varphi$ \emph{तार्किकरीत्या
  निष्पन्न होते}, म्हणजे $\Gamma \Entails \varphi$, तेव्हा आणि
  केवळ तेव्हाच, जेव्हा प्रत्येक प्रतिमान
  $\mModel{M} = \tuple{W, R, V}$ आणि प्रत्येक $w \in W$
  साठी, सर्व $\psi \in \Gamma$ बाबत $\mSat{M}{\psi}[w]$
  असेल तर $\mSat{M}{\varphi}[w]$ असते. $\Gamma$ मध्ये
  एकच सूत्र~$\psi$ असेल, तर $\psi \Entails \varphi$ लिहितो.
\end{defn}

\begin{ex}
  एका सूत्र पासून दुसरे तार्किकरीत्या निष्पन्न होते हे
  दाखवण्यासाठी $\mSat{M}{}[w]$ ची व्याख्या वापरून सर्व
  प्रतिमानांबद्दल विचार करावा लागतो. उदाहरणार्थ,
  $p \lif \Diamond p \Entails \Box\lnot p \lif \lnot p$
  दाखवण्यासाठी असा युक्तिवाद करता येईल: प्रतिमान
  $\mModel{M} = \tuple{W, R,
    V}$ आणि $w \in W$ घ्या, आणि $\mSat{M}{p \lif \Diamond
    p}[w]$ असे समजा. $\mSat{M}{\Box\lnot p \lif \lnot
    p}[w]$ दाखवायचे आहे. तसे नाही असे समजा. मग
  $\mSat{M}{\Box\lnot p}[w]$ आणि $\mSat/{M}{\lnot p}[w]$.
  $\mSat/{M}{\lnot p}[w]$ असल्यामुळे $\mSat{M}{
    p}[w]$. गृहीतकावरून $\mSat{M}{p \lif \Diamond p}[w]$;
  म्हणून $\mSat{M}{\Diamond p}[w]$.
  $\mSat{M}{\Diamond
    p}[w]$ च्या व्याख्येवरून $Rww'$ असलेले असे $w'$
  आहे की $\mSat{M}{p}[w']$. पण $\mSat{M}{\Box \lnot p}[w]$
  देखील असल्यामुळे $\mSat{M}{\lnot p}[w']$; हा विरोधाभास.

  सूत्र~$\psi$ पासून दुसरे~$\varphi$ तार्किकरीत्या
  निष्पन्न होत नाही हे दाखवण्यासाठी प्रतिदृष्टांत द्यावा
  लागतो: असे प्रतिमान $\mModel{M} = \tuple{W, R,
    V}$ की त्यातील एखाद्या जगात~$w \in W$ येथे
  $\mSat{M}{\psi}[w]$, पण $\mSat/{M}{\varphi}[w]$.
  $p \lif \Diamond p \Entails/
  \Box p \lif p$ हे दाखवू. \ref{nml:syn:ent:fig:counterex} मधील
  प्रतिमान पाहा. येथे $\mSat{M}{\Diamond p}[w_1]$
  आणि म्हणून $\mSat{M}{p \lif
    \Diamond p}[w_1]$. मात्र $\mSat{M}{\Box p}[w_1]$
  असूनही $\mSat/{M}{p}[w_1]$ असल्यामुळे
  $\mSat/{M}{\Box p \lif p}[w_1]$.
  \begin{figure}
  \begin{center}
    \begin{tikzpicture}[modal]
      \node[world] (w1) [label=right:\mFalse{p}]{$w_1$};
      \node[world] (w2) [label=right:\mTrue{p}, above left=of w1]{$w_2$};
      \node[world] (w3) [label=right:\mTrue{p}, above right=of w1] {$w_3$};
      \draw[->] (w1) to (w2);
      \draw[->] (w1) to (w3);
    \end{tikzpicture}
  \end{center}
  \caption{$p \lif \Diamond p
  \Entails \Box p \lif p$ याचा प्रतिदृष्टांत.}
  \label{nml:syn:ent:fig:counterex}
  \end{figure}

  अनेकदा अगदी साधे प्रतिदृष्टांत पुरेसे असतात.
  $W' = \{w\}$, $R' = \emptyset$, आणि $V'(p) =
  \emptyset$ असलेले प्रतिमान $\mModel{M'} =
  \tuple{W', R', V'}$ हाही प्रतिदृष्टांत आहे:
  $\mSat/{M'}{p}[w]$ असल्यामुळे
  $\mSat{M'}{p \lif \Diamond p}[w]$. $w$ पासून एकही
  जग प्राप्य नसल्यामुळे $\mSat{M'}{\Box p}[w]$; म्हणून
  $\mSat/{M'}{\Box p
    \lif p}[w]$.
\end{ex}

\begin{prob}
  $\Box (\varphi \land \psi) \Entails \Box \varphi$ हे दाखवा.
\end{prob}

\begin{prob}
  $\Box (p \lif q) \Entails/ p \lif \Box q$ आणि $p \lif \Box
  q \Entails/\Box (p \lif q)$ हे दाखवा.
\end{prob}



